\section{Introduction}
\label{introduction}

Modern deep learning (DL) programs increasingly depend on libraries such as Transformers~\cite{transformers} and Megatron~\cite{megatron}, which streamline model development, training, and deployment~\cite{zeng2024survey,alfadel2023empirical,wu2023understanding,lyu2023chronos,kula2018developers,samarinunderstanding,chen2023toward,pepe2024taxonomy,li2024knowbug,jia2021symptoms,otoom2019automated}. 
These libraries encapsulate critical functionalities such as tokenization, optimizer scheduling, distributed communication, and memory management, enabling developers to build large-scale models efficiently.

However, as these libraries grow in complexity and scale, they unavoidably introduce software defects. 
Bugs in DL libraries may not always manifest as program crashes; instead, they often degrade performance silently, introduce numerical inconsistencies, or corrupt training dynamics under certain conditions~\cite{chen2023toward,yu2025towards}. 
For example, a subtle error in tensor parallelism coordination or configuration parsing may lead to incorrect gradient updates or inefficient resource usage, ultimately impairing model convergence.
In practice, such defects are often discussed in GitHub pull requests (PRs)~\cite{pr} or commits~\cite{commit}, where developers propose fixes based on user feedback or internal audits~\cite{xiao2024generative,herbold2022fine}.
Thus, we can analyze PRs and commits to assess whether a defect impacts downstream client programs.
However, this is highly non-trivial in practice: PRs may reference internal APIs, abstract trigger conditions, or rely on complex runtime configurations. 
Meanwhile, downstream users may not closely monitor upstream libraries and are frequently unaware of whether such defects could impact their own programs.

Although keeping dependencies up-to-date is a common practice for maintaining software reliability, frequent dependency upgrades are impractical in DL environments, as each update may trigger costly retraining and break compatibility~\cite{2024Understanding,2023Toward,2024Bug,10.5555/3767901.3767919}.
Moreover, many defects remain silent across releases or are fixed with minimal documentation.
Release notes are typically coarse (e.g., fix stability) and omit critical details such as affected modules, configurations, or triggering conditions.
For instance, commit \texttt{1674ce3}~\cite{1674ce3} in \textit{Megatron} silently skewed multi-user training for months before detection, illustrating how such fixes can propagate unnoticed.

Recent efforts have leveraged large language models (LLMs) to assist with code summarization and PR triage~\cite{liu2019automatic}. 
However, these approaches typically treat commits or patches in isolation, without modeling the downstream impact on external client programs.
Determining whether a DL library defect affects a specific client program involves several unique challenges:

\noindent\textbf{Challenge 1: Noisy and informal defect representations}.
Defect-related PRs or commits often contain unstructured descriptions, partial fixes, or discussion noise. 
The root cause and triggering conditions are rarely formalized, making the automated extraction of defect patterns difficult.

\noindent\textbf{Challenge 2: Deep semantic gap between library changes and user-visible behavior}.
Many patches fix low-level internal logic (e.g., CUDA graphs, communication buffers), while user programs interact only with high-level APIs. 
Mapping these low-level changes to user-facing risk factors (e.g., configuration parameters, method calls) requires semantic reasoning.

\noindent\textbf{Challenge 3: Semantic alignment between defect conditions and client usage}. 
Determining whether a client program satisfies the defect's triggering conditions requires understanding not only the patched library logic, but also the usage context in the user code—including configurations and function calls. 

To address these challenges, we propose \appname, an agent coordination framework for automated defect and impact analysis in deep learning libraries. 
\appname assists developers in determining whether a defect in an upstream DL library (e.g., Transformers) has concrete effects on their downstream programs. 
Unlike prior approaches that rely on static dependency resolution or flat summarization, \appname leverages a modular agent architecture to reason across both the library side and the client side, enabling context-aware, condition-sensitive impact assessment.
By integrating domain-specific coordination rules, agent self-adaptation, and AST-based structural validation, \appname turns generic multi-agent reasoning into a practical and verifiable workflow for DL defect analysis.
\appname consists of four agents: (1) \textbf{Miner Agent} identifies whether a commit or pull request addresses a real defect and extracts structured defect metadata, including bug background, trigger conditions, and affected components.
(2) \textbf{Code Diff Agent} analyzes the patch semantics, locating modified methods and summarizing the underlying fix logic.
(3) \textbf{Orchestrator Agent} synthesizes the results from upstream agents into a defect pattern, capturing risk-relevant method calls and parameter combinations from the perspective of external users.
(4) \textbf{Impact Analyzer Agent} performs contextual analysis on the client program code to determine whether it satisfies the triggering conditions, guided by static analysis and domain-specific validation rules.

To mitigate the long-context limitations of LLMs, \appname incorporates a progressive context augmentation mechanism that dynamically expands or compresses the input prompt to ensure information completeness without exceeding token limits.
Furthermore, to mitigate hallucinations in LLM-generated outputs, we design a post-processing validation module that uses AST-based static analysis to verify the existence of referenced parameters and methods within the actual client program. 
% If hallucinations are detected, the agent is prompted to regenerate its output, up to a fixed retry threshold.

We evaluate \textbf{\appname} on 157 PRs and 70 commits collected from two widely used DL libraries: \texttt{Transformers} and \texttt{Megatron}.
Our study focuses on two core tasks: (1) whether \appname can accurately generate structured \textit{defect patterns}, and (2) whether it can determine if a downstream client program is semantically affected by a given defect.
On the PR dataset, \appname achieves a defect classification F1-score of 95\%, and over 70\% of extracted defect pattern fields—such as bug background and trigger conditions—are rated fully accurate by domain experts.
In downstream impact analysis across 122 real-world client programs, \appname achieves an F1-score of 85\%, substantially outperforming three strong baselines.
We further evaluate \appname on internal commits of \texttt{Megatron}, where it successfully identifies 12 real-world client programs affected by silent defects, as confirmed by developers.

In summary, the main contributions of this paper include:
\begin{itemize}[leftmargin=*]
\item We present \textbf{\appname}, a novel agent coordination framework for fine-grained defect and impact analysis in deep learning libraries. 
\item We implement a practical system that integrates LLM-based agents with static analysis and domain-specific inference rules. 
The replication package is available at~\cite{depradar}.
\item We conduct an extensive evaluation of 227 library updates (157 PRs and 70 commits) from \texttt{Transformers} and \texttt{Megatron}. 
Results show that \appname outperforms three strong baselines in both defect modeling (F1 = 95\%) and client impact analysis (F1 = 85\%), and identifies real affected clients.
\end{itemize}

